\begin{abstract}
% Every number below is in Results: Tables 2-3, Sec. 4.4-4.6.
Test-time augmentation (TTA) is a cheap way to raise mean segmentation scores, but averages can hide lesions a model stops finding. We evaluated geometric TTA with 2, 4 and 8 dihedral views as a pre-deployment question on four models: three 2D polyp segmentation models of different quality, trained on CVC-ClinicDB with sequence-level cross-validation (15 checkpoints each), and a 3D SegResNet for enhancing tumour on BraTS 2021. The 2D models were tested without retraining on three external polyp collections (1576 images) and on 1009 polyp-free colonoscopy frames. TTA raised Dice for every model on every dataset with lesions, and raised NSD and lowered HD95 among cases both arms segmented. Whether it lost lesions under shift depended on how a miss was counted, and the disagreement holds with the checkpoint as the unit of analysis. For SwinUNETR, a model that failed by design, images with no overlap rose in all 15 checkpoints on all three collections while missed images at an intersection over union (IoU) of 0.10 moved on none, so within a collection, on the same predictions, one criterion says significantly harmed and the other says nothing. For UNet the two corrected findings point opposite ways in different collections: more images left empty on CVC-ColonDB in all 15 checkpoints and fewer missed on ETIS-Larib in all 15 (medians $+10$ of 380 and $-11$ of 196, Holm $p=0.002$ each). Requiring IoU 0.10 we could not detect a rise in missed images for any model, at any threshold up to 0.50. UNet's false alarms on polyp-free frames fell from 70.8\% to 56.5\%, and for empty outputs TTA was close to the curve traced by changing the threshold of a single view, within a difference we could not resolve. TTA moved the operating point with no significant net loss of lesions a model placed on target, but within intervals wide enough to hide one: the widest admits about 20 more missed images per 1000. TTA should be validated on external and lesion-free data, with misses counted by an explicit localisation criterion and reported next to empty predictions and false alarms.
\end{abstract}

\begin{keywords}
test-time augmentation, segmentation, reliability, external validation, missed detection, uncertainty
\end{keywords}
